\section{Finite framed groupoids and the small genera}
\label{sec:B}

We use $E(G)$ for the internal edges, $I(G)$ for the individually
labelled legs, and $F(G)$ for all flags.  Thus the manuscript's
$\widehat A$ is $J(G):=E(G)\amalg I(G)$, and its
$\widehat{\vec A}$ is $F(G)$.  The map
\[
 q:F(G)\longrightarrow J(G)
\]
has two-element fibres over internal edges and one-element fibres over
legs.  All products in the first subsection are on the \emph{pinned}
stratum $M_G=\prod_v M_{g(v),F_v}$; we introduce the action of
$\operatorname{Aut}(G)$ explicitly when we pass to mixed groupoids.
The $S$ in cote 146 retains the tangent directions at the legs as well
as the smoothing directions at the nodes.  This convention is visible
both in the first quotient on scan 100 and in scans 103--110.
The primary manuscript sources are \citet{B_Grothendieck146}
and \citet{B_Grothendieck148}; scan numbers below refer to their
PDF pages, with manuscript folios stated where visually checked.

\subsection{The contracted product, including its finite sign action}

\begin{result}{B.1}{Proved}{The framed torsor and its two quotients}
For a pinned stable graph $G$, put
\[
 P_G=\prod_{j\in J(G),\,M_G} U(L_j),\qquad
 L_i=T_{x_i}C\quad(i\in I(G)),\qquad
 L_e=T_{x_h}C_{v(h)}\otimes T_{x_{\bar h}}C_{v(\bar h)}.
\]
Here $U(L)=L^{\times}/\mathbb R_{>0}$ is the oriented unit-ray circle
bundle; no Hermitian metric is required.  We write
$P_G^E=\prod_{e\in E(G)}U(L_e)$.  Choose a finite set $B_G$ in the
pinned stratum, a nonempty finite set of rays above each member of
$B_G$, and saturate the rays under
$\mu_G=\{\pm1\}^{J(G)}$.  Denote the resulting finite set by $D_G$.
When a graph action is required, take $B_G$ invariant and also
saturate $D_G$ under the finite group $\operatorname{Aut}(G)$.
For vertices of genus zero and valence three or four, the groupoid
\[
 \mathcal S_G:=\Pi_1(P_G,D_G)
\]
is a $\mu_G$-free $\mathbb Z^{J(G)}$-groupoid, and its internal lattice
acts freely on every nonempty Hom-set.  Its geometric quotients are
\[
 \mathcal S_G/(\mathbb Z^{I(G)},\mu_2^{I(G)})
       \simeq\Pi_1(P_G^E,\operatorname{im}D_G),
 \qquad
 \mathcal S_G/(\mathbb Z^{J(G)},\mu_G)
       \simeq\Pi_1(M_G,B_G).
 \tag{B.1}
\]
The first groupoid has the node-framed link group of Part A.  The
second has automorphism group $\prod_{v:\,|F_v|=4}F_2$.  After choices
of trivializations, the automorphism group of $\mathcal S_G$ is
\[
 \mathbb Z^{J(G)}\times\prod_{v:\,|F_v|=4}F_2.
 \tag{B.2}
\]
The direct-product splitting in (B.2) is not specified as
$\operatorname{Aut}(G)$-equivariant.
\end{result}
\begin{proof}
For a trivalent rational component the moduli space is a point.  For
four flags its moduli space is $\mathbb P^1\setminus\{0,1,\infty\}$,
which retracts to a graph of rank two.  Each tangent-ray bundle on
that graph is trivial: construct a section on a maximal tree, and
extend it around each remaining edge using the connectedness of
$S^1$.  Each $L_j$ above is a tensor product of pullbacks of such
bundles.  Consequently their product admits a section, and $P_G$ is
isomorphic, after this choice, to
$M_G\times(S^1)^{J(G)}$.  This proves (B.2) and proves injectivity of
the full lattice without an appeal to the notes' claim of freeness.

The finite sign action is multiplication of nonzero rays.  A sign
fixing a ray in every coordinate has all coordinates $+1$, so its
action on $D_G$ is free.  For paths with prescribed endpoints, the
projection to the base has fibres that are torsors under the winding
lattice.  Dividing by the leg-coordinate winding lattice and
identifying the corresponding sign-orbits of endpoints therefore
gives exactly the first groupoid in (B.1); dividing by all coordinates
gives the second.  In particular, the quotient notation includes
both the identification of endpoint lifts and the killing of their
vertical loops.  It does not mean the action groupoid of the finite
sign group.

For completeness, the contraction from flags to edges is the map of
torus bundles
\[
 \prod_{h\in F(G)}U(T_{x_h}C_{v(h)})\longrightarrow P_G,
 \qquad (u_h)_h\longmapsto
 \bigl((u_hu_{\bar h})_{e=\{h,\bar h\}},(u_i)_{i\in I}\bigr).
 \tag{B.3}
\]
Its map on integer lattices sends both $e_h,e_{\bar h}$ to $e_e$,
and its map on signs multiplies the two signs.  Its kernel is the
connected torus with one coordinate per edge, embedded by
$u\mapsto(u,u^{-1})$.  Killing its full fundamental groupoid,
including the paths between its different sign lifts, is the
contracted product in scan 100.  To check the assertion on each
fibre, lift the circles to $\mathbb R^{F(G)}$ and
$\mathbb R^{J(G)}$: the same coordinate-sum map is surjective, and
its kernel is generated by $e_h-e_{\bar h}$.  Two lifted paths have
the same image precisely when their difference is a kernel path.
This proves the universal quotient statement on arrows and
composition.  Applying the construction over the product base
proves (B.3) on fundamental groupoids.  Graph automorphisms permute
flags and edges and commute with (B.3), so these constructions carry
the required graph action before taking mixed quotients.
\end{proof}

\begin{result}{B.2}{Proved}{The presentation actually written in scans 115--117}
For a trivalent rational graph set $r=|J(G)|$.  The finite model of
cote 146, scans 114--117 (folios 113--116), has the $2^r$ objects of
the sign torsor $D_G=\prod_{j\in J(G)}\varepsilon_j$.  It has
$r2^r$ positive half-turn generators
\[
 u_{j,x}:x\longrightarrow\sigma_jx
 \quad(j\in J(G),\ x\in D_G).
\]
Its relations are the source-compatible commuting squares for two
different coordinates.  The commutation of the full-turn loops
$v_{j,x}=u_{j,\sigma_jx}u_{j,x}$ written on scan 117 follows from
these squares.  Equivalently,
\[
 \operatorname{Hom}(x,y)=
 \{m\in\mathbb Z^r:m\bmod2=y-x\},
 \tag{B.4}
\]
where addition gives composition and $m$ records displacement in
units of a half turn.  The loop lattice is $2\mathbb Z^r$, identified
with $\mathbb Z^r$ by $2e_j\mapsto e_j$.
\end{result}
\begin{proof}
Choose a member of the sign torsor only to write its members as
binary vectors.  The universal covering cubulation of the torus has
vertex set $\mathbb Z^r$ and oriented coordinate edges; quotient by
$2\mathbb Z^r$ gives the displayed vertices and generators.  A
sequence of coordinate edges has a displacement $m$.  Commuting
squares interchange adjacent steps in different coordinates, and
backtracking cancels inverse steps.  Sorting coordinate steps and
then cancelling proves that two paths with the same endpoints have
the same class exactly when their displacement vectors agree.
This proves (B.4), with no unproved relation among winding numbers.
The two paths consisting of two steps in coordinate $i$ and two in
coordinate $j$ are interchanged by four commuting squares, proving
the stated redundancy of scan 117's second relation.

A half-turn is therefore an arrow between two distinct chosen
framings.  Its square is a full twist; imposing order two on this
arrow would kill a nonzero member of the meridian lattice.  The
formula also proves the cocycle on scan 116: the discrepancy between
two chosen binary half-turns and their chosen product is one full
turn in precisely the coordinates where both signs are negative.
\end{proof}

\begin{result}{B.3}{Refuted}{The proposed equivalence before forgetting leg framings}
The proposed equivalence $S_0\mathcal T_G\simeq S_A\Pi M_G$ without
quotienting the leg framings is false, already for the graph with
one trivalent rational vertex and three labelled legs.  The exact
comparison is the first functor in (B.1), whose kernel at every
object is $\mathbb Z^{I(G)}$.
\end{result}
\begin{proof}
For this graph $M_G$ is a point and $E(G)=\varnothing$.  Formula
(B.4) gives automorphism group $\mathbb Z^3$ for the scan 115 model;
the early two-object model of scans 3--8 has the same automorphism
group.  The groupoid $S_A\Pi M_G$ has trivial automorphism group.
The forgetful functor sends $(m_1,m_2,m_3)$ to zero and is not
faithful.  No equivalence between the two groupoids exists because
an equivalence induces a bijection on each endomorphism set.  The
quotient by this displayed kernel is the terminal groupoid, which
is the required node-only model.  For general $G$ the same argument
in each fibre gives the comparison (B.1).
\end{proof}

\begin{result}{B.4}{Proved}{The genus-zero tree index and the chart counts}
The category of reduced stable trees with a fixed individually
labelled leaf set, and specified edge contractions fixing those
labels, is equivalent to a finite poset.  On four and five leaves,
the numbers of trees by number of internal edges are respectively
$(1,3)$ and $(1,10,15)$.  For the disjoint union of their trivalent
charts, the scan 116 presentations have the following counts;
these counts precede any global identification of chart objects.
\[
\begin{array}{c|r|r|r|r|r}
 n&\text{charts}&r&\text{framed objects}&\text{half-turn generators}
       &\text{node-only objects / generators}\\\hline
 4&3&5&96&480&6\,/\,6\\
 5&15&7&1920&13440&60\,/\,120
\end{array}
\tag{B.5}
\]
\end{result}
\begin{proof}
Every internal edge separates the labelled leaves into two nonempty
sets, each of cardinality at least two.  Distinct edges give distinct
bipartitions.  The collection of these bipartitions reconstructs the
tree: the vertices are the compatible regions determined by the
edge cuts, and incidence is given by changing one region across one
cut.  Thus an automorphism fixing every leaf fixes every edge and
then every vertex.  A contraction removes precisely its contracted
edge bipartitions.  Its source and target therefore determine the
contracted set uniquely.  This proves both absence of nontrivial
automorphisms and uniqueness of a contraction, the two ingredients
needed for a poset.  A stable tree on $n$ leaves has at most $n-3$
internal edges, so there are finitely many collections of the finite
set of leaf bipartitions.

For four leaves there are three bipartitions into two pairs.  For
five leaves there are ten bipartitions of type $2+3$; two such cuts
are compatible exactly when their two-element parts are disjoint.
There are $\binom52\binom32/2=15$ such pairs.  A trivalent tree has
$n-3$ internal edges, hence $r=2n-3$ framed coordinates.  Apply B.2
and sum over the three or fifteen charts.  The program
\texttt{code/b\_finite\_checks.py} independently enumerates stable
multigraphs, fixes all labels, checks stability, and obtains all
entries of (B.5).  This also specifies the precise finite counting
convention; adding further marked basepoints changes a presentation's
object and generator counts.
\end{proof}

\begin{result}{B.5}{Refuted}{Freeness does not follow from an unspecified finite colimit}
The implication that a finite colimit of groupoids with a free
$\mu_2$-action on objects has a free action on objects is false.
Consequently that implication alone cannot justify the sentence
following the first formula of cote 146, scan 100.
\end{result}
\begin{proof}
Let $C$ be the indiscrete groupoid on the two objects $0,1$, and let
$\mu_2$ interchange them.  The two equivariant functors
$\operatorname{id},\sigma:C\rightrightarrows C$ have coequalizer
with one object.  Its inherited $\mu_2$-action fixes that object.
More explicitly, if $u:0\to1$ is the unique arrow, the coequalizer
identifies $u$ with $u^{-1}$, and has automorphism group $C_2$.
The action on objects is therefore no longer free.  The program
checks the two-object coequalizer.  In B.1 freeness was proved from
the action on actual nonzero rays; it was not deduced from the
existence of a colimit.  The present counterexample addresses the
inference about arbitrary finite colimits, without substituting
an index category for the unfinished index on scan 100.
\end{proof}

\begin{result}{B.6}{Proved}{An explicit obstruction for the boundary-incidence interpretation}
Let $L$ be the boundary of a sufficiently small regular
neighbourhood of the boundary divisor in $\overline M_{0,5}$.  The
colimit of the fundamental groupoids of its circle-bundle pieces
and all their torus intersections presents $\Pi_1(L)$.  The inclusion
$L\to M_{0,5}$ does not induce an isomorphism of fundamental groups.
Thus this specified boundary-incidence colimit cannot be the
missing global presentation in scan 100.
\end{result}
\begin{proof}
There are ten rational boundary components, each meeting three
others, and fifteen intersections.  Their incidence graph has the
ten two-element subsets of a five-element set as vertices and
joins disjoint subsets.  It is the Petersen graph.  Removing small
discs around the intersections leaves ten pairs of pants.  Above
each is a circle bundle, necessarily trivial since a pair of pants
has the homotopy type of a graph.  Thus each vertex group is
$\mathbb Z\times F_2$, each intersection group is $\mathbb Z^2$,
and both edge inclusions are injective: one generator is the circle
fibre, and the other is an infinite-order peripheral loop in the
pair of pants.  Van Kampen with basepoints on all intersections
proves the asserted colimit.

The universal cover obtained by developing this graph of spaces is
a tree of contractible universal covers of these pieces, glued
along contractible universal covers of the tori.  Its finite
subtrees are contractible by successive gluing along contractible
subspaces, and every compact sphere lies in a finite subtree.
Thus $L$ is aspherical.  It is a connected closed oriented
three-manifold, so
\[
 H_3(\pi_1L;\mathbb Z)=H_3(L;\mathbb Z)=\mathbb Z.
\]
On the other hand, the map
$M_{0,5}\to M_{0,4}$ forgetting the fifth point has fibre a sphere
minus four points.  Both base and fibre are graphs up to homotopy,
of ranks two and three.  A section can be given in affine
coordinates $(x,y)$ by $y=2+|x|$.
The resulting extension is $F_3\rtimes F_2$.  Its classifying space
is a graph of mapping cylinders of homotopy equivalences of a
rank-three graph, one cylinder for each free base generator.
Its universal cover is a tree of contractible graph covers and
mapping cylinders, so it is a $K(\pi,1)$ of dimension two.  Hence its third homology is zero.  An isomorphism induced by
$L\to M_{0,5}$ would identify these third homology groups, a
contradiction.  This proof specifies the actual comparison map and
the invariant that prevents it from being an equivalence; it makes
no assertion that the manuscript had defined this index category.
\end{proof}

\begin{result}{B.7}{Known}{The precise relation to the Lego groupoid}
Let $\mathcal L(\Sigma)$ be the marking two-complex of
Bakalov--Kirillov with their moves and relations, and let $\Gamma$
be the mapping class group with the same orientation and boundary
conventions.  The groupoid with objects the markings and arrows
\[
 M\longrightarrow M':\quad (g,[p]),\qquad
 g\in\Gamma,\quad p:gM\leadsto M',
\]
is equivalent to $B\Gamma$ by $(g,[p])\mapsto g$.  With fixed labels
and capped boundary circles it gives the open groupoid
$B\Gamma_{0,n}$ in Part A.  The capping functor kills exactly the
central lattice of boundary twists.
\end{result}
\begin{proof}
Theorem 4.9 of \citet{B_BakalovKirillov} states connectedness and
simple connectedness of $\mathcal L(\Sigma)$ in genus zero; its
proof in \S6 reduces to the sphere mapping-class presentation in
Proposition 6.6 and the commuting lifts of cutting moves.  Therefore
for any $g,M,M'$ there is exactly one relative-endpoint path class
$p$.  The displayed map is bijective on every Hom-set and
essentially surjective.  Composition is
\[
 (h,[q])\circ(g,[p])=(hg,[q*(hp)]),
\]
where the path first traverses $hp$.  This proves the equivalence
with its exact composition map.  Capping sends a homeomorphism
fixed on every boundary to a homeomorphism fixing the new labelled
punctures.  A puncture-fixing isotopy can be adjusted near each
puncture to preserve a small circle; its residual winding is an
integer twist in that collar.  Conversely those collar twists
become isotopic to the identity after capping.  This proves both
surjectivity and the stated kernel.  It is the same quotient as the
leg quotient in B.1.  These maps give a comparison of the groupoids.
They do not assign a word to each arrow in the different, unfinished
finite diagram of cote 148; those words require a chosen set of
geometric paths for that diagram.
\end{proof}

\begin{result}{B.8}{Open}{The complete finite colimit and the generator dictionary}
The index category, transition functors, and relations giving the
literal general finite colimit of cote 146, scan 100 have not been
reconstructed here.  The trivalent fibre presentation B.2, the
four-valent product and its two quotients B.1, and the finite tree
index B.4 are proved.  A generator-by-generator comparison between
that unspecified colimit and the Lego presentation is not proved.
\end{result}
The early four-point model is now reconstructed in B.16: the
thirty objects, forty-eight essential arrows, eighteen relations
and all ninety-six expanded generator words are specified there.
B.17 gives its free-sign induction to four leg framings and tracks
the dual scalar convention.  These local computations do not
specify the general diagram indexed in scan 100.  Scans 121--125
select square and tetrahedral fibres in a later construction;
scan 124 postpones relations and scan 125 ends mid-sentence.
The colimit inference B.5, the boundary-incidence map B.6, and
the equivariant marking functor B.7 have their exact scopes stated.
The remaining reconstruction is the transition functor on every
index arrow, together with a complete proof that its colimit
relations agree with those of the geometric fundamental groupoid.
This row records that undecided item, not an alternative completion
of Part B(b).

\subsection{A finite colimit whose maps are fully specified}

\begin{result}{B.8a}{Proved}{Finite presentation yields a finite colimit of elementary framed groupoids}
For each pinned stable graph, its framed fundamental groupoid admits
a finite-colimit presentation in the category of groupoids with its
specified central lattice and finite sign action.  One may take the
elementary automorphism groups, after changing the central lattice,
to be $\mathbb Z^r$ and $\mathbb Z^r\times F_2$, the groups of the
three- and four-point rational pieces.  The arrows in this colimit
are the word maps specified below.  This proves existence of a
finite algebraic presentation of the asserted kind.  It does not
identify these word maps with the unfinished tree-indexed diagram
of scan 100.
\end{result}
\begin{proof}
Put $\pi=\mathbb Z^r$ for the framed coordinates and let $H$ be the
framed group.  The quotient $\Gamma=H/\pi$ is finitely presented.
For genus zero this follows from the braid presentation and capping;
for positive genus it follows from
\citet[Theorems 3.1 and 3.2]{B_LabruereParis}, followed by the finite
Schreier rewriting for the subgroup fixing all labels.  A product
of vertex groups is finitely presented; gluing two boundary
coordinates imposes one central relation per edge.  Thus the same
assertion holds for every pinned stable graph.

Write $\Gamma=\langle x_1,\ldots,x_m\mid R_1,\ldots,R_s\rangle$
and choose lifts $\widetilde x_i\in H$.  Each lifted relator equals
a uniquely determined element $c_j\in\pi$.  Then
\[
 H=\left(\pi\times F(x_1,\ldots,x_m)\right)
       /\langle\!\langle R_jc_j^{-1}:1\leq j\leq s\rangle\!\rangle.
 \tag{B.6}
\]
To verify injectivity of this presentation into $H$, take a word
mapping to the identity.  Its image in $\Gamma$ is a product of
conjugates of the $R_j^{\pm1}$.  Replace these relators by their
specified central values using (B.6).  What remains lies in $\pi$
and maps to zero in its injective copy in $H$; hence it is zero.
Surjectivity follows by lifting a word in the $x_i$ and adjusting
its residual element of $\pi$.

Let $E=\tfrac12\mathbb Z^r$, with the inclusion
$\pi\subset E$, and form $H_0=H\amalg_{\pi}E$ with $E$ central.
There is a homomorphism $\chi:H_0\to E/\pi=\mu_2^r$ that kills the
chosen $\widetilde x_i$ and restricts to the quotient on $E$.
The translation groupoid of $H_0$ acting on $\mu_2^r$ by $\chi$ is
the framed groupoid on one full sign orbit; its vertex group is
$H$.  Formula (B.6), now with $E$ in place of $\pi$, gives its
finite-colimit construction explicitly: for each relator use the
parallel maps
\[
 E\times F(t)\rightrightarrows E\times F(x_1,\ldots,x_m),
 \qquad t\mapsto R_jc_j^{-1},\quad t\mapsto1,
 \tag{B.7}
\]
with both maps the identity on $E$.  Their coequalizer imposes
exactly (B.6).  The group $E\times F_m$ is a finite amalgamated
coproduct over $E$ of its one-generator free factors.  The
one-generator factor itself is the coequalizer of the identity
and the map $t\mapsto t,u\mapsto1$ on $E\times F(t,u)$.
Consequently every group appearing as a piece can be replaced by
$E$ or $E\times F_2$ in a finite diagram.  All maps preserve
$\chi$ and induce functors that are the identity on the sign-object
set.  Colimits on arrows therefore give the translation groupoid
of $H_0$, with its free sign action.  The proof verifies the
particular diagram's freeness and avoids the false general
inference of B.5.  Additional finitely many basepoints are added
by chosen connecting isomorphisms, which also constitute a finite
diagram.
\end{proof}

\subsection{Explicit presentations in types $(1,1)$ and $(1,2)$}

In the next statements $a,b,c$ are positive Dehn twists on a
three-chain: consecutive curves intersect once and the first and
third are disjoint.  With one marked point we use only $a,b$.
The choices consist of an oriented reference surface, these curves,
and paths from the chosen framed objects to this reference.  The
following homomorphisms are genuine homomorphisms for those choices.
Changing a reference path conjugates the corresponding
homomorphism by its endpoint path; we do not call the resulting
based isomorphisms canonical.

\begin{result}{B.9}{Proved}{The two open groupoids and their leg-framed extensions}
For complex curves with individually labelled markings,
\[
 \Gamma_{1,1}=\langle a,b\mid aba=bab,\ (ab)^6=1\rangle
                 \cong\operatorname{SL}_2(\mathbb Z),
 \tag{B.8}
\]
\[
 \Gamma_{1,2}=\langle a,b,c\mid aba=bab,\ bcb=cbc,
               \ ac=ca,\ (abc)^4=1\rangle.
 \tag{B.9}
\]
Thus $\Pi_1M_{1,n}$ is equivalent to the one-object groupoid on the
displayed group.  The groups of the unit tangent-ray bundles at
all marked points are respectively
\[
 H_{1,1}=\langle a,b\mid aba=bab\rangle,
       \qquad z_1=(ab)^6,
 \tag{B.10}
\]
and
\[
 H_{1,2}=\left\langle a,b,c,z_1,z_2\ \middle|\
 \begin{array}{l}
 aba=bab,\ bcb=cbc,\ ac=ca,\\
 z_1,z_2\text{ central},\ (abc)^4=z_1z_2
 \end{array}\right\rangle
 \cong B_4\times\mathbb Z.
 \tag{B.11}
\]
Here $z_i$ is the positive full turn at the $i$th marked point.
The capping maps are the quotients $z_i\mapsto1$.
\end{result}
\begin{proof}
For two labelled points $p,q$ on a smooth genus-one curve, the
complete linear system $|p+q|$ gives a double cover of a projective
line, with $p,q$ over infinity and four distinct finite branch
points.  An affine coordinate $x$ on the line is unique up to
$x\mapsto\alpha x+\beta$.  Ordering $p,q$ fixes the sign of $y$
at infinity in a monic equation
\[
 y^2=\prod_{i=1}^4(x-r_i).
\]
Every isomorphism of the pointed curves induces the affine map of
the quotient and has exactly one lift preserving the ordering of
$p,q$.  Conversely that lift gives the pointed-curve isomorphism.
These constructions apply in families, giving the analytic-stack
equivalence
\[
 M_{1,2}\simeq[\operatorname{UConf}_4(\mathbb C)/
                        \operatorname{Aff}(\mathbb C)].
\]
The unordered configuration space has braid group $B_4$, with its
usual adjacent-crossing generators and relations.  The translation
subgroup of the affine group is contractible; the positive loop in
its dilation subgroup rotates all four roots once, giving the
central braid $(abc)^4$.  The homotopy exact sequence for this
quotient stack gives exactly (B.9).  The full rotation injects into
$B_4$, since its image under exponent sum is $12$.

At either point over infinity, $u=1/x$ is a local coordinate.
Under dilation it is multiplied by $\alpha^{-1}$.  Keeping two
nonzero tangent vectors replaces the preceding quotient by the
quotient of $\operatorname{UConf}_4(\mathbb C)\times(\mathbb C^*)^2$
with these two weights.  The affine loop therefore imposes exactly
$(abc)^4z_1^{-1}z_2^{-1}=1$.  This proves (B.11), including the
identification of the two central generators and the capping map.
Eliminating $z_2$ exhibits $B_4\times\mathbb Z$.

For one labelled point the double cover from $|2p|$ has infinity
ramified and three finite branch points.  In a monic cubic equation
the allowed scaling is $x\mapsto\lambda^2x+\beta$,
$y\mapsto\lambda^3y$.  A positive loop of $\lambda$ rotates the
three roots twice, giving $(ab)^6$.  The local parameter $x/y$ at
$p$ has weight $\lambda^{-1}$.  The same calculation proves
(B.8) and (B.10).  Mapping $a,b$ to
\[
 \begin{pmatrix}1&1\\0&1\end{pmatrix},\qquad
 \begin{pmatrix}1&0\\-1&1\end{pmatrix}
\]
identifies (B.8) with $\operatorname{SL}_2(\mathbb Z)$: the braid
presentation becomes the standard amalgam
$C_4*_{C_2}C_6$ on $aba$ and $ab$.
The same formulas are obtained by specializing
\citet[Theorem 3.1 and Proposition 2.12]{B_LabruereParis} to
one and two boundary components.  In the two-boundary case their
Artin diagram is the chain $a,b,c$ together with one central
boundary generator; their chain relation identifies the other
boundary generator with $(abc)^4z_1^{-1}$.
\end{proof}

\begin{result}{B.10}{Proved}{All boundary graph groups in these two types}
For $(1,1)$ there is one boundary graph: a rational vertex with one
loop and one labelled leg.  Its pinned stratum is a point and its
mixed stratum is $BC_2$.  Put $h=(ab)^3$.  Its fully framed mixed
link group is
\[
 \langle a,h\mid[a,h]=1\rangle,\qquad z_1=h^2.
 \tag{B.12}
\]
Its node-only mixed link group is
$\langle a,h\mid[a,h]=1,h^2=1\rangle$.

For $(1,2)$ name the four boundary graph types as follows:
$N$ is a rational vertex with a loop and both legs;
$T$ is an elliptic vertex joined to a rational tail carrying both
legs; $C$ is a two-edge rational cycle with one labelled leg on
each vertex; and $D$ is a rational loop vertex joined to a rational
tail carrying both legs.  Write
\[
 s=aba,\quad\Delta=abcaba,\quad
 r=s^{-1}\Delta s,\quad q=(abc)^2,\quad h=(ab)^3,
 \quad d=h^2.
 \tag{B.13}
\]
The pinned stratum groups are $F_2$, $\operatorname{SL}_2(\mathbb Z)$,
$1$, $1$ for $N,T,C,D$, respectively.  The mixed stratum groups are
$C_2*\mathbb Z$, $\operatorname{SL}_2(\mathbb Z)$, $C_2$, $C_2$.
The following are presentations of their fully framed mixed link
groups; in every row $z_1,z_2$ are central:
\[
\begin{array}{c|l|l}
 \text{graph}&\text{generators and additional relations}
                   &\text{node twists}\\\hline
 N&a,c,r,z_1,z_2;\ [a,c]=[a,r]=1,\ r^2=z_1z_2&a\\
 T&a,b,z_1,z_2;\ aba=bab&d=(ab)^6\\
 C&a,c,q,z_1,z_2;\ [a,c]=1,\ qaq^{-1}=c,\ q^2=z_1z_2&a,c\\
 D&a,h,z_1,z_2;\ [a,h]=1&a,d=h^2
\end{array}
\tag{B.14}
\]
The pinned full-framing subgroups are the kernels of the graph
permutation: in rows $N,C,D$ the respective quotient maps send
$r,q,h$ to the nonidentity member of $C_2$, and all lattice
coordinates to the identity.  The graph group in row $T$ is
trivial.  The node-only link groups are obtained by $z_1=z_2=1$.
The mixed stratum groups are obtained by also killing the listed
node twists.
\end{result}
\begin{proof}
The graph enumeration is finite.  For one marking it gives the
open elliptic graph and the one-loop graph.  For two markings it
gives the open graph and exactly $N,T,C,D$.  The program enumerates
all connected multigraphs with at most $3g-3+n$ edges, assigns vertex
genera with total genus $g$, assigns the individual legs, and checks
$2g(v)-2+\operatorname{val}(v)>0$.  It obtains graph automorphism
orders $2,1,2,2$ for $N,T,C,D$, as asserted.

Cutting the node of $N$ leaves a four-point sphere.  Pin the two
external markings at $0,\infty$ and the two node branches at $1,z$.
Exchanging the branches acts by $z\mapsto z^{-1}$.  Its quotient has
coordinate $w=z+z^{-1}$, two missing points $2,\infty$, and a single
orbifold point of order two at $-2$.  Its orbifold fundamental group
is therefore $\langle r,c\mid r^2=1\rangle=C_2*\mathbb Z$.
The subgroup fixing the branches is free on $c,rcr^{-1}$.  The
node circle supplies a commuting generator $a$.  At the order-two
configuration, the involution of the four-point sphere is
$x\mapsto-x$: its derivatives at the two external markings are
$-1$, and the product of its derivatives at the two branches is
$+1$.  A framed lift consequently has square $z_1z_2$, with no
node-winding factor.  This proves row $N$ by the presentation of a
central extension of $C_2*\mathbb Z$; there is no additional
relator on its free generator $c$.

Cutting the separating node of $T$ leaves a one-point elliptic
component and a rigid rational component.  Keeping the node
framing converts the elliptic factor to $B_3$ by (B.10); the two
external leg framings supply $z_1,z_2$.  Its boundary twist is
$(ab)^6$.  This proves row $T$.

For $C$ both rational components are rigid.  The node and leg
lattices have rank four.  The graph involution interchanges the two
nodes, and its lift $q$ interchanges $a,c$; its square is $z_1z_2$
by (B.11).  The index-two subgroup generated by
$a,c,z_1,z_2$ is free abelian: its generators are independent twists
about the four boundary or cutting circles.  Cutting into annuli
and pairs of pants proves that a product of these twists is
trivial only when every exponent is zero.  Thus the extension by
$C_2$ has precisely row $C$'s presentation.

For $D$, the involution flips the branches of the rational loop.
After smoothing the loop it is the elliptic half-turn $h$ in the
one-holed torus, whose square is the separating-node twist $d$.
It commutes with $a$ and both external twists.  The subgroup on
$a,h^2,z_1,z_2$ is again the independent rank-four lattice, so the
displayed extension has no further relation.  The same argument
with one external marking gives (B.12).  All passages to pinned
and node-only groups are the specified kernel and quotient maps;
the finite graph actions have not been discarded.
\end{proof}

\begin{result}{B.11}{Proved}{Generization on every boundary flag in the two small genera}
The generization maps from the groups (B.12)--(B.14) to the open
framed groups are the inclusions specified by the words (B.13).
In $(1,2)$ the two nontrivial two-step flags, in addition to the
cycle flags, have maps
\[
 D\longrightarrow T:\quad a\mapsto a,\ h\mapsto(ab)^3,
 \qquad
 D\longrightarrow N:\quad a\mapsto a,\ h\mapsto ac^{-1}r,
 \tag{B.15}
\]
with $z_i\mapsto z_i$.  For $C\to N$ retaining the node $a$, the
arrow-stabilizer group is the subgroup
$\langle a,c,z_1,z_2\rangle$ of row $C$, and its map to row $N$ is
the identity on these named generators.  The other contraction is
its conjugate by $q$.  These maps compose to the maximal
generization maps.  After the lattice quotients in B.10 they give
the corresponding pinned and mixed groupoid spans of Part A.
\end{result}
\begin{proof}
A positive rotation of a plumbing parameter is the positive Dehn
twist about its vanishing circle.  Therefore each lattice
coordinate has the stated image.  The remaining generators are
homeomorphisms of the cut components extended over the annuli;
this gives the indicated homomorphisms into the stabilizers of the
remaining curves.  The braid relations give
\[
 \Delta^2=(abc)^4,\quad rar^{-1}=a,\quad
 q a q^{-1}=c,\quad q c q^{-1}=a.
\]
They also give the exact word identity needed in (B.15):
\[
 \begin{split}
 ac^{-1}r
 &=ac^{-1}a^{-1}b^{-1}a^{-1}(abcaba)(aba)\\
 &=(aba)^2=(ab)^3=h.
 \end{split}
\]
Thus both routes from $D$ to the open group agree as genuine based
homomorphisms.  The same is true of the cycle maps on their
arrow-stabilizers.  The element $q$ exchanges the two cycle edges,
so it does not preserve the choice of the one edge being smoothed;
this is exactly why it is absent from that arrow-stabilizer.  There
is no map from all of row $C$ to that chosen branch of row $N$
asserted here.  Including its second contraction gives the other
branch and the required graph equivariance.

For verification the program \texttt{b\_braid\_checks.py} evaluates
both sides of all the displayed braid identities in the faithful
Artin action on a free group.  Its twelve tests all pass.  The
proof of faithfulness is not needed to infer the homomorphisms,
which follow from the geometric construction; the computation is
an independent exact check of the words used in that construction.
The complete lists of boundary graphs and contractions show that
no further flags occur in these two types.
\end{proof}

\subsection{The finite five-point diagram, checked against its original image}

\begin{result}{B.12}{Proved}{The nine vertex orbits and the corrected sixteen arrow orbits}
Let $K$ be a fixed five-element set and let $\mathfrak S_K$ act by
relabeling.  The vertex sets of cote 148, scan 31, are the following
nine transitive $\mathfrak S_K$-sets.  The symbols in the first column
are the manuscript's symbols; subscripts are stabilizer orders.
\[
\begin{array}{c|c|r}
\text{type}&\text{stabilizer in }S_5&\text{number}\\ \hline
\mathrm I_{10}&D_5&12\\
\mathrm I_6&S_3\text{ acting on }3+2\text{ by its sign}&20\\
\mathrm I_4&C_4&30\\
\mathrm {II}_2&C_2&60\\
\mathrm {II}'_2&C_2&60\\
\mathrm {II}_1&1&120\\
\mathrm {III}_2&C_2&60\\
\mathrm {III}_1&1&120\\
\mathrm {III}'_2&C_2&60
\end{array}
\tag{B.16}
\]
In every $C_2$ row the nonidentity element is a double
transposition.  Thus the number of vertices is $542$, in nine
orbits.  The fourteen named arrow families have the following
cardinalities and orbit decompositions:
\[
\begin{array}{c|c|r|r}
\text{family}&\text{source and target}&\text{arrows}&S_5\text{-orbits}\\ \hline
 f_{B,s}&\mathrm I_6\to\mathrm I_4&60&1\\
 f_{\vec Q,a}&\mathrm {II}'_2\to\mathrm I_4&120&1\\
 \tau_{\pi,s}&\mathrm {II}_2\to\mathrm {II}'_2&60&1\\
 \tau_r^0&\mathrm {II}_1\to\mathrm {II}_1&120&1\\
 \gamma_{\pi,s}&\mathrm {II}_2\to\mathrm I_{10}&60&1\\
 \gamma''_{\pi,s}&\mathrm {II}'_2\to\mathrm I_6&60&1\\
 \gamma_r^0&\mathrm {II}_1\to\mathrm I_6&120&1\\
 \tau_{\pi,\rho}&\mathrm {III}_2\to\mathrm {III}_1&120&1\\
 \gamma_{\pi,\rho}&\mathrm {III}_2\to\mathrm {II}_2&120&1\\
 \tau_r^1&\mathrm {III}_1\to\mathrm {III}'_2&120&1\\
 \gamma_r^1&\mathrm {III}_1\to\mathrm {II}'_2&120&1\\
 \gamma_r^{1'}&\mathrm {III}_1\to\mathrm {II}_1&360&3\\
 \gamma_r^2&\mathrm {III}'_2\to\mathrm {II}_2&120&1\\
 \gamma_r^{2'}&\mathrm {III}'_2\to\mathrm I_4&60&1
\end{array}
\tag{B.17}
\]
Consequently $1620$ arrows means sixteen orbits, with orbit-size
histogram $5\cdot60+11\cdot120$.  The reported version ``fourteen
orbits'' is false; fourteen counts the named families in (B.17).
The original scan already strikes out the word \emph{orbites} next
to $14$, so this correction is not attributed to an unqualified
retained assertion of the author.
\end{result}
\begin{proof}
An unoriented pentagonal order on $K$ has dihedral stabilizer of
order ten, giving $120/10=12$ objects of type $\mathrm I_{10}$.
A bitetrahedral structure is a partition $K=J\amalg\varepsilon$,
$|J|=3$, together with a bijection from $\varepsilon$ to the two
orientations of $J$.  There are $\binom53\cdot2=20$.  Its stabilizer
permutes $J$ arbitrarily and acts on $\varepsilon$ by the sign of
that permutation.  An oriented square on $K\setminus\{P\}$ has
cyclic stabilizer of order four, and there are $5\cdot6=30$ such
objects.  A pentagon with a distinguished vertex has stabilizer
the reflection fixing that vertex, a double transposition; there
are $12\cdot5=60$.  This describes the two $\mathrm {II}$ types
with subscript two and $\mathrm {III}_2$.  A pinned structure of
$\mathrm {II}_1$ has data
\[
 K=J\amalg\varepsilon,\quad |J|=3,\quad
 \omega\in\operatorname{Or}(J),\quad r\in J\times\varepsilon.
\]
There are $10\cdot2\cdot6=120$, and fixing all these data fixes
each label.  The fully pinned $\mathrm {III}_1$ structures have
the same trivial stabilizer.  The square with a specified
codiagonal frame at $\mathrm {III}'_2$ has a remaining half-turn,
and hence $120/2=60$ objects.  These constructions prove (B.16)
as identifications of $S_5$-sets, not just as a sum of integers.

For the arrows, the first family has parameters $(B,s)$ with
$s\in J(B)$; it has $20\cdot3=60$ members and double-transposition
stabilizer.  The second has an oriented square and one of its four
specified flags, hence $30\cdot4=120$ members and trivial
stabilizer.  The families parameterized by $(\pi,s)$ have the
pentagon-with-vertex stabilizer.  The families parameterized by a
complete frame $r$, or by $(\pi,\rho)$ where $\rho$ is a pentagonal
frame, have trivial stabilizer.  This verifies every single-orbit
row of (B.17).  The three choices of gommage at a fully pinned
$\mathrm {III}_1$ vertex give three arrows at each of its 120
members.  An equivariant orbit of arrows over a free transitive
source set contains exactly one arrow over every source: a
stabilizer of an arrow is contained in the trivial source
stabilizer.  Therefore these $360$ arrows split into precisely
three free orbits.  Adding gives the claimed sixteen orbits.

Two scan checks are essential here.  The long bent arrow labeled
$\gamma_r^2$ starts at $\mathrm {III}'_2$, not at
$\mathrm {III}_1$ as in the supplied TikZ transcription.  Its
end-valencies are $2,2$, agreeing with $120=2\cdot60=2\cdot60$
on scan 33.  The diagonal $\gamma_r^{1'}$ on scan 31 has valency
three at both of its 120-member ends.  Scan 33 retains
$3\cdot120$ on its left but adds the arithmetically incompatible
$3\cdot40=3\cdot40$ on its right.  Both the valencies and the
retained subtotal $1020$ force $360$ for that row.  We record the
incompatible reading rather than silently replacing its $40$.

The program \texttt{b\_finite\_checks.py} enumerates all $120$
permutations and all the displayed cosets.  It produces each
vertex and each member of each named arrow family, independently
reproducing (B.16)--(B.17).  Its endpoint maps are explicitly typed
coset maps $gH\mapsto(gH_s,gkH_t)$ with
$H\subset H_s\cap kH_tk^{-1}$.  The generated
\texttt{b\_diagram.json} states its precise scope: it is a
realization of these endpoint valencies, and does not assign the
geometric homotopy classes of the manuscript's paths.
\end{proof}

\begin{result}{B.13}{Proved}{The relations required by the five-point diagram}
The underlying geometric graph of the diagram in (B.16)--(B.17)
is connected.  The free groupoid on its displayed arrows has
isotropy group $F_{1079}$.  The fundamental group of $M_{0,5}$ has
abelianization $\mathbb Z^5$.  Therefore a presentation of
$\Pi_1M_{0,5}$ on exactly those vertices and arrows must impose
relations whose images in the abelianization of $F_{1079}$ span
a subgroup of rank $1074$.
\end{result}
\begin{proof}
The bipartite subgraph between $\mathrm I_6$ and $\mathrm I_4$
has a concrete coset description.  Label $K=\{0,1,2,3,4\}$ and take
\[
 H_6=\langle(012),(12)(34)\rangle,
 \qquad H_4=\langle(1324)\rangle,
 \qquad H_2=\langle(12)(34)\rangle.
\]
Its edges are $S_5/H_2$, with endpoints $gH_6,gH_4$.
The distinguished square vertex is $0$, and its diagonals are
$\{1,2\},\{3,4\}$; thus this is exactly the incidence obtained by
choosing the distinguished member of the three-element equator
of a bitetrahedron.  The group $\langle H_6,H_4\rangle$ is $S_5$.
For an elementary verification, conjugating $x=(012)$ by
$(1324)$ gives $y=(034)$.  Conjugating $x$ by $y,y^2$ gives
$(312),(412)$.  Thus all $(i12)$, $i\in\{0,3,4\}$ occur,
and products $(i12)(j12)^{-1}=(i1j)$ give all three-cycles
containing $1$.  The identity $(1ab)(1bc)=(abc)$ then gives
every three-cycle, hence $A_5$.  The four-cycle is odd.  Enumeration in the attached program also
checks the order $120$.  A path in this bipartite coset graph is a
product of elements of $H_6,H_4$, proving connectedness.  Every
other vertex type has a displayed arrow toward this subgraph or
toward a type already connected to it.  In particular
$\mathrm I_{10}$ connects through $\mathrm {II}_2$ and its
$\tau$-bijection with $\mathrm {II}'_2$.  Each incidence has
positive valency at every endpoint, so this proves connectedness
of the full graph.

Choose a spanning tree of the graph.  Successively contracting
its 541 edges gives a groupoid equivalence to the one-object free
group on $1620-541=1079$ generators, one for each remaining edge.
The comparison sends an arrow $e:v\to w$ to the loop obtained by
following the tree path from a fixed base vertex to $v$, then $e$,
then the inverse tree path from the base vertex to $w$.  This also
specifies the equivalence on every arrow.

Normalize three of the five markings to $0,1,\infty$.  Then
\[
 M_{0,5}=\{(x,y)\in\mathbb C^2:
 x,y\ne0,1,\quad x\ne y\}.
\]
Forgetting $y$ is a locally trivial bundle with base
$\mathbb C\setminus\{0,1\}$ and fibre
$\mathbb C\setminus\{0,1,x\}$.  A section is
$y=2+|x|$, since it avoids $0,1,x$.  It gives
$\pi_1(M_{0,5})=F_3\rtimes F_2$.  The base pure braids act trivially
on the abelianization of the fibre: each labelled puncture
meridian is sent to a conjugate of itself.  Abelianizing this
semidirect product therefore gives $\mathbb Z^3\oplus\mathbb Z^2$.
The rank statement follows from the exact abelianization of a
group presentation.  It requires at least $1074$ independent
abelianized relators; it does not purport to list a complete
nonabelian relation system.
\end{proof}

\begin{result}{B.14}{Proved}{The pentagon as an explicit relation among pants moves}
For a five-holed sphere with a fixed cyclic order and a fixed
outer boundary, let $T_0,\ldots,T_4$ be the five binary planar
pants trees.  Choose the standard compatible marking on each.
Write $A_i:T_i\to T_{i+1}$ for the elementary change of one
internal cut, with indices modulo five.  In the marking groupoid
of \citet[Sections 4.5, 4.7 and 4.10]{B_BakalovKirillov}, each
$A_i$ is the specified composite of a cut-removal move and a
cut-insertion move, and
\[
 A_4A_3A_2A_1A_0=1_{T_0}. \tag{B.18}
\]
Under the equivariant comparison of B.7 this is a genuine
relation in the mapping-class-group groupoid.  Under capping
boundary circles it remains a relation in $\Pi_1M_{0,5}$.
\end{result}
\begin{proof}
Let $U$ be the uncut marking with the chosen cyclic order and
outer distinguished edge.  For each $T_i$, remove its two cuts
and denote the resulting path to $U$ by $p_i$.  The generalized
$F$-move is defined in the cited Section 4.5 with a specified
distinguished edge at its target; hence all these maps have
specified objects and target marking.  The associativity-of-cuts
relation, equation (4.7) of that paper, says that the two orders
of removal give the same $p_i$.

Adjacent trees share one cut.  Let $N_i$ be the one-cut marking
obtained from either by removing the nonshared cut, and write
$f_i:T_i\to N_i$ and $f'_i:T_{i+1}\to N_i$ for the respective
cut-removal moves.  Set $A_i=(f'_i)^{-1}f_i$.  Removing the remaining
cut from $N_i$ gives a path $b_i:N_i\to U$.  Associativity gives
$p_i=b_if_i$ and $p_{i+1}=b_if'_i$, and hence
\[
 p_{i+1}A_i=p_i,
 \qquad A_i=p_{i+1}^{-1}p_i.
\]
Multiplying these five equations cancels the adjacent $p_i$'s
and proves (B.18).  Every move in $A_i$ is supported on the
four-holed subsurface where the pants decomposition changes.
The relation uses the common uncut five-holed marking; it is the
specific dimension-two compatibility absent from a presentation
consisting only of unrelated one-dimensional pieces.

More explicitly, the equivariant groupoid has arrows $(g,[p])$
with $p:gM\to M'$.  The word in (B.18) maps to
$(1,[A_4\cdots A_0])=(1,[1])$ by the homotopy just written.
Applying the boundary-twist quotient functor in B.7 preserves
this equality.  Thus both comparison maps and the path homotopy
are specified; no identification of an elementary move with an
unbased homeomorphism is used.
\end{proof}

\begin{result}{B.15}{Open}{The remaining literal presentation tasks in B(b) and B(d)}
The computations B.12--B.13 settle the finite cardinalities,
the orbit correction, and the amount of relation data required
in the five-point diagram.  B.14 supplies the explicit
four-point-move pentagon in the established marking presentation.
A complete list of relators on the actual $1620$ geometric arrows
of cote 148, together with the requested expression of every
one of them by relative-dimension-at-most-one arrows, has not
been obtained here.  In particular the endpoint coset maps in
\texttt{b\_diagram.json} do not determine those path classes and
cannot be used as such a list.

The missing comparison is a word assignment to the geometric
$\gamma_r^{1'}$ and $f_{B,s}$ families, compatible with the
specified quarter-turn, third-turn and sixth-turn paths and the
framed central lattices.  The early four-point arrows
$\widetilde a_r,\widetilde b_r,\widetilde c_r,\widetilde d_r$
are now assigned words in B.16, with the inverse and changed source
in (B.21) stated explicitly.  That calculation proves its local
presentation; it does not supply the five-point family dictionary
or the literal general finite-index diagram.  The chart counts in
B.4 and the induced-model counts in B.17 refer to their stated
presentations.  The precise global colimit and the five-point word
assignment remain the unproved steps in Part B; the abstract
finite-colimit construction B.8a does not supply them.
\end{result}

\subsection{A direct reconstruction of the early four-point model}

\begin{result}{B.16}{Proved}{Thirty objects, forty-eight arrows, and the four-point word reduction}
The arrangement model underlying cote 146, scans 20, 27, 31 and
33 (folios 19, 26, 30 and 32), admits the following complete typed
presentation.  It has thirty objects, forty-eight essential
arrows and eighteen defining relations.  The expanded generator
table of scan 27, including its full-turn and half-turn auxiliary
arrows, has ninety-six arrows.  Its isotropy group is
\[
 \langle A,F,t\mid[t,A]=[t,F]=1\rangle=F(A,F)\times\mathbb Z(t).
 \tag{B.19}
\]
The central $t$ is one positive complex-scalar turn in the
arrangement complement.  Quotienting by it gives the four-point
moduli group $F_2$.  The explicit words below give a comparison
functor on every generator, including the auxiliary ones.
\end{result}

\begin{proof}
Here are the data without uncertain subscripts.  Indices $k$ lie
in $\mathbb Z/6$, and $\epsilon$ is $+1$ or $-1$.  In a Euclidean
plane $V$ place three lines at angles $0,\pi/3,2\pi/3$ and write
$J$ for its positive quarter-turn.  Put
\[
 X=V_{\mathbb C}\setminus\bigcup_{j=0}^2(L_j)_{\mathbb C},
 \qquad Q_k=(\cos((k+\tfrac12)\pi/3),
                   \sin((k+\tfrac12)\pi/3)).
\]
Choose $0<\eta<\pi/12$.  In the chamber containing $Q_k$, choose
$R_{k,+}$ at angle $(k+1)\pi/3-\eta$, and $R_{k,-}$ at angle
$k\pi/3+\eta$, on the unit circle.  Finally put
\[
 P_{k,\epsilon}=Q_k-i\epsilon JQ_k.
\]
There are six $Q$'s, twelve $R$'s and twelve $P$'s.  On each of
the two complex eigenlines of $J$, the six $P$'s are consecutive
scalar sixth-roots, exactly the pole circles and the three object
families of the original diagrams.  Changing the sufficiently
small $\eta$ supplies the canonical path transport within the
same chamber and does not change this presentation.

The four families of arrows, each with twelve members, are
\[
\begin{array}{c|c}
 b_{k,\epsilon}&Q_k\to R_{k,\epsilon}\\
 c_{k,\epsilon}&P_{k,\epsilon}\to Q_k\\
 a_{k,\epsilon}&R_{k,\epsilon}\to R_{k+\epsilon,-\epsilon}\\
 d_{k,\epsilon}&P_{k,\epsilon}\to P_{k+\epsilon,\epsilon}.
\end{array}
\tag{B.20}
\]
Take $b$ to be the straight segment in the real chamber, take
$c(s)=Q_k-i\epsilon(1-s)JQ_k$, and take
$d(s)=e^{\pi i s/3}P_{k,\epsilon}$.  The identity
$J P_{k,\epsilon}=i\epsilon P_{k,\epsilon}$ gives precisely the
stated endpoint of $d$.  For $a$, let $u,v$ be its source and target
and set
\[
 a(s)=\tfrac12(u+v)+e^{\pi i s}\tfrac12(u-v).
\]
This is a positive half-turn in the complex normal coordinate to
the intervening line.  None of these paths meets an excluded
complex line: for $b,c$ the appropriate nonzero real linear
functional has constant sign in its real part; for $d$ its real
part rotates through two adjacent chambers and its imaginary
part prevents a zero at their common wall.  For $a$ the two other
line functionals have modulus of their constant term strictly
larger than that of their rotating term, by the bound on $\eta$.

To match the manuscript's arrow orientations, let
$r(k,\epsilon)$ denote the incidence of the intervening wall with
the chamber containing $Q_k$.  Then $b,c,d$ in (B.20) are its
$\widetilde b,\widetilde c,\widetilde d$ with these incidences,
whereas our positive crossing is
\[
 a_{k,\epsilon}
   =\widetilde a_{r(k+\epsilon,-\epsilon)}^{-1}.
 \tag{B.21}
\]
This orientation is forced by its scan 31 mixed equation
$\widetilde a\,\widetilde b\,\widetilde c\,\widetilde d
 =\widetilde b\,\widetilde c$: the $\widetilde a$ there traverses
back from the terminal chamber of $\widetilde d$.  Thus we specify
an inverse and its changed source, rather than identify two
oppositely directed arrows.

The twelve mixed relations are
\[
 c_{k+\epsilon,\epsilon}d_{k,\epsilon}
 =b_{k+\epsilon,-\epsilon}^{-1}a_{k,\epsilon}
       b_{k,\epsilon}c_{k,\epsilon}.                 \tag{B.22}
\]
Define the paths from $Q_k$ to $Q_{k+3}$
\[
 H_{k,\epsilon}=
 c_{k+3,\epsilon}
 d_{k+2\epsilon,\epsilon}d_{k+\epsilon,\epsilon}d_{k,\epsilon}
 c_{k,\epsilon}^{-1}.
\]
The six further relations are
\[
 H_{k,+}=H_{k,-}\quad(k\in\mathbb Z/6).             \tag{B.23}
\]
They are the typed half-turn equality represented by the
manuscript's uncertain products of $\widetilde A$'s in scans
31 and 33.  Three consecutive sixth-turns give a half-turn;
six give a full turn.  These path lengths determine the invariant
meaning without assigning a new paleographic reading to an
uncertain exponent.

We first verify these equations geometrically.  Both sides of
(B.22), after removing their common initial $c$, lie in the domain
where the real parts of the two noncrossed line functionals have
their prescribed signs.  In coordinates normal and tangent to
the crossed wall this domain has the form
\[
 \{(z,w)\in\mathbb C^*\times\mathbb C:
          \operatorname{Re}w>c\,|\operatorname{Re}z|\}
\]
after a real linear rescaling, for some $c>0$.  Its projection to
$\mathbb C^*$ has convex fibres, so relative path classes are
classified by winding in $z$.  Both paths have the same endpoints
and a positive argument change of $\pi$ in $z$.  This proves
(B.22).  For (B.23), the square
\[
 (u,s)\longmapsto e^{\pi i s}(Q_k-iu\epsilon JQ_k),
 \qquad 0\leq u,s\leq1,
\]
lies in $X$: every excluded real line functional has a nonzero
real part before the scalar factor.  Its boundary identifies
$H_{k,\epsilon}$ with the scalar half-turn
$\mu_{Q_k}(s)=e^{\pi i s}Q_k$.  Both signs give the same path
class, as required.

These are all defining relations.  Eliminate the twelve $R$
objects using the arrows $b$, and the twelve $P$ objects using
$c$.  Define the resulting wall crossings
\[
 e_{k,\epsilon}
   =b_{k+\epsilon,-\epsilon}^{-1}a_{k,\epsilon}b_{k,\epsilon}.
\]
The mixed equations eliminate all twelve $d$'s.  What remains has
six chamber objects, twelve directed wall crossings, and the six
relations equating the two length-three galleries to the
opposite chamber.  This is the $A_2$ chamber presentation.  One
may check completeness directly by identifying the three lines
with $x_i=x_j$ in $\{x_1+x_2+x_3=0\}$: adding the centroid gives
an isomorphism $\mathbb C\times X\cong\operatorname{Conf}_3
(\mathbb C)$.  The six chamber objects are its six real orderings,
and its adjacent normal half-turns are the standard adjacent
braid crossings.  Generic braid diagrams generate all paths;
a generic homotopy between diagrams has only cancellation of
opposite crossings or the three-strand move
$\sigma_1\sigma_2\sigma_1=\sigma_2\sigma_1\sigma_2$.
Lifting that relation to each ordering gives exactly (B.23).
Equivalently these are the six vertices, twelve edges and six
two-cells of the arrangement complex described in
\citet[Section 4.1]{B_SalvettiSettepanella}.  This verifies that
the path functor from (B.20)--(B.23) to $\Pi_1(X)$ is fully
faithful and essentially surjective.

For a completely algebraic word check, put
$x_k=e_{k,+}$ and $y_k=e_{k,-}$, and contract the spanning tree
$x_0,\ldots,x_4$.  Set $x_5=t$, $y_0=A$, $y_5=F$.  The six
remaining relations solve as follows:
\[
\begin{array}{c|rrrrrr}
k&0&1&2&3&4&5\\ \hline
x_k&1&1&1&1&1&t\\
y_k&A&A^{-1}F^{-1}&F&At&A^{-1}F^{-1}&F.
\end{array}                                                   \tag{B.24}
\]
They leave precisely $[t,A]=[t,F]=1$.  For example, the first
three equations give $y_4=y_1=A^{-1}F^{-1}$ and $y_2=F$; the
fourth gives $y_3=At$; the last two impose the displayed
commutators.  Conversely substitution into all six equations
verifies them, so no relator has been lost.  Sending every $b,c$
to $1$ and every $a_{k,\epsilon},d_{k,\epsilon}$ to the value
of $e_{k,\epsilon}$ in (B.24) gives the announced functor to
(B.19).  Every eliminated arrow belongs to a specified spanning
tree or is solved by a displayed relation, proving that this
functor is an equivalence.

The scalar full turn is the product of six successive $d$'s,
and maps to $t$ for either sign.  Every $\lambda$ in scan 27
therefore maps to $t$.  The $\mu_{Q_k}$ maps to $1$ for
$k=0,1,2$ and to $t$ for $k=3,4,5$; its target is $Q_{k+3}$,
so this does not identify a half-turn with an identity on its
original object.  The $\mu_R$ words are obtained by transport
along the specified $b$ arrows.  This accounts for all thirty
$\lambda$'s and all eighteen $\mu_Q,\mu_R$'s in that table.
The $\lambda$-transport relations follow from centrality of $t$.

Finally the affine normalization
\[
 (z_1,z_2,z_3)\longmapsto
 \bigl(z_1,z_2-z_1,(z_3-z_1)/(z_2-z_1)\bigr)
\]
identifies the ordered configuration space with
$\mathbb C\times\mathbb C^*\times
 (\mathbb P^1\setminus\{0,1,\infty\})$.
A scalar full turn is the positive generator of its
$\mathbb C^*$ factor, and the projective quotient kills exactly
$t$.  This proves the claimed comparison with $M_{0,4}$.
The program \texttt{b\_four\_point.py} lists all objects,
all forty-eight arrows and all eighteen typed relations, checks
both endpoints and the words (B.24), and checks all twelve pole
full turns.  Its output also verifies the expanded count $96$.
\end{proof}

\begin{result}{B.17}{Proved}{Four leg framings and the dual scalar convention}
The free-sign induction of the thirty-object model B.16 from
$(\mathbb Z,\mu_2)$ to $(\mathbb Z^4,\mu_2^4)$ along the diagonal
has $240$ objects.  A product presentation with three additional
coordinate half-turn circles has $1104$ essential generators,
or $1488$ when retaining the auxiliary generators of scan 27.
Its automorphism group is $\mathbb Z^4\times F_2$, and its
quotient by all four leg framings is $\Pi_1M_{0,4}$.
These counts name this particular induced version of the early
model; finite equivalent presentations need not have the same
object count.

The two scalar conventions in cote 146, scans 18 and 50--53,
are related by an explicit dualization functor.  On the full
winding lattice it is multiplication by $-1$, and on the finite
sign group it is the identity.  Hence the positive scalar $t$
of B.16 is the inverse of the positive principal-bundle turn
when the manuscript calls the scalar action ``opposite'' on
scan 53.  The corresponding lattice pushout is respectively
$t\mapsto(1,1,1,1)$ or $t\mapsto(-1,-1,-1,-1)$, with this
comparison map specified.
\end{result}
\begin{proof}
The diagonal lattice is primitive.  Complete its generator
$w=(1,1,1,1)$ to the basis $w,e_2,e_3,e_4$ of $\mathbb Z^4$.
Modulo two, the same basis splits the sign extension into the
old antipodal action and three additional independent sign
coordinates.  Thus the induced finite groupoid is the product
of B.16 with three copies of the two-object half-turn circle
of B.2.  The basis is a stated choice, not a permutation-invariant
splitting.  This gives $30\cdot2^3=240$ objects and
\[
 48\cdot8+30\cdot(3\cdot8)=1104
\]
generators.  Using the unreduced $96$ in place of $48$ gives
$1488$.  The product relations consist of the lifted eighteen
relations, the commuting squares in the three circle coordinates,
and the squares commuting an old arrow with a new coordinate
half-turn.  The product universal property proves completeness.
Its central lattice is the pushout of $\mathbb Z(t)$ along the
stated diagonal, so its group is $\mathbb Z^4\times F_2$.
Killing that lattice and identifying its sign lifts gives the
projective quotient in B.16, hence $\Pi_1M_{0,4}$.

For the convention comparison, let $L$ be a complex line bundle.
For each nonzero vector $v$ there is exactly one complex-linear
functional $\ell_v$ with $\ell_v(v)=1$.  Thus
$v\mapsto\ell_v$ is a canonical continuous map
$L^{\times}\to(L^{\vee})^{\times}$.  It satisfies
$\ell_{\alpha v}=\alpha^{-1}\ell_v$, descends to the ray circles,
and sends a positive scalar loop to a negative scalar loop.
It commutes with multiplication by $-1$ and is natural under
linear isomorphisms of line bundles.  This defines the functor
on paths as well as on objects, proves its lattice and sign
formulas, and tracks the reversal that scan 53 states.
Applying this functor before the diagonal pushout gives precisely
the second displayed lattice map.  No unproved identification
between a line and its dual, with unchanged positive inertia,
is made.
\end{proof}
